\uuid{7mB4}
\titre{ EMV loi de Laplace – minimisation par la médiane empirique }

\niveau{L3}
\module{ Statistique inférentielle }
\chapitre{Maximum de vraisemblance}
\sousChapitre{}
\theme{Statistique}
\auteur{Maxime Nguyen}
\datecreate{ 2026-03-31}
\organisation{}
\difficulte{4}

\contenu{

\texte{ 
Soit $X_1,\dots,X_n$ des variables aléatoires i.i.d. de densité commune
\[
f_\theta(x)=\frac12 e^{-|x-\theta|},\qquad \theta\in\mathbb{R}.
\]
On note $X_{(1)}\leq X_{(2)}\leq \cdots \leq X_{(n)}$ les statistiques d’ordre associées à l’échantillon.
}

\begin{enumerate}

\item   \question{Écrire la vraisemblance $L(\theta)$ et la log-vraisemblance $\ell(\theta)$. Montrer que maximiser $\ell(\theta)$ est équivalent à minimiser
\[
g(\theta)=\sum_{i=1}^n |X_i-\theta|.
\]}
\indication{}
\reponse{
La vraisemblance est
\[
L(\theta)=\prod_{i=1}^n \frac12 e^{-|X_i-\theta|}
=\left(\frac12\right)^n \exp\left(-\sum_{i=1}^n |X_i-\theta|\right).
\]
En prenant le logarithme :
\[
\ell(\theta)=\ln L(\theta)
=-n\ln 2-\sum_{i=1}^n |X_i-\theta|
=-n\ln 2-g(\theta).
\]
Puisque $-n\ln 2$ est une constante, maximiser $\ell(\theta)$ revient à minimiser
\[
g(\theta)=\sum_{i=1}^n |X_i-\theta|.
\]
}

\item   \question{Montrer que la fonction $\theta\mapsto |x-\theta|$ est convexe pour tout $x\in\mathbb{R}$, et en déduire que $g$ est convexe.}
\indication{}
\reponse{
Pour tout $x$ fixé, la fonction
\[
h:\theta\mapsto |x-\theta|
\]
est convexe. En effet, pour tous $\theta_1,\theta_2\in\mathbb{R}$ et tout $\lambda\in[0,1]$,
\[
h(\lambda \theta_1+(1-\lambda)\theta_2)
=
|x-\lambda\theta_1-(1-\lambda)\theta_2|
\]
\[
=
|\lambda(x-\theta_1)+(1-\lambda)(x-\theta_2)|
\leq \lambda |x-\theta_1|+(1-\lambda)|x-\theta_2|
\]
\[
=\lambda h(\theta_1)+(1-\lambda)h(\theta_2).
\]
Donc $h$ est convexe.

Comme
\[
g(\theta)=\sum_{i=1}^n |X_i-\theta|
\]
est une somme finie de fonctions convexes, $g$ est elle-même convexe.
}

\item   \question{Montrer que $g$ est dérivable sur chaque intervalle ouvert $(X_{(k)},X_{(k+1)})$ pour $k=0,1,\dots,n$ (avec la convention $X_{(0)}=-\infty$ et $X_{(n+1)}=+\infty$), et calculer $g'(\theta)$ sur chacun de ces intervalles en fonction de $k$.}
\indication{}
\reponse{
Sur l’intervalle $(X_{(k)},X_{(k+1)})$, aucune observation ne coïncide avec $\theta$, donc chaque terme $|X_i-\theta|$ est dérivable en $\theta$.

Plus précisément,
\[
\frac{d}{d\theta}|X_i-\theta|=
\begin{cases}
1 & \text{si } X_i<\theta,\\
-1 & \text{si } X_i>\theta.
\end{cases}
\]
Sur $(X_{(k)},X_{(k+1)})$, exactement $k$ observations vérifient $X_i<\theta$ et $n-k$ vérifient $X_i>\theta$. Donc
\[
g'(\theta)=k\cdot 1 +(n-k)\cdot (-1)=k-(n-k)=2k-n.
\]
}

\item   \question{Étudier le signe de $g'(\theta)$ sur chaque intervalle et en déduire les variations de $g$. Montrer que $g$ est décroissante puis croissante, et identifier l’intervalle (ou le point) où $g$ atteint son minimum.}
\indication{}
\reponse{
Sur l’intervalle $(X_{(k)},X_{(k+1)})$, on a
\[
g'(\theta)=2k-n.
\]
Donc :
\begin{itemize}
\item si $k<\dfrac n2$, alors $g'(\theta)<0$ : $g$ est strictement décroissante ;
\item si $k>\dfrac n2$, alors $g'(\theta)>0$ : $g$ est strictement croissante ;
\item si $k=\dfrac n2$ (ce qui n’est possible que si $n$ est pair), alors $g'(\theta)=0$ : $g$ est constante sur l’intervalle correspondant.
\end{itemize}

Ainsi, $g$ est décroissante puis croissante. Son minimum est atteint au passage entre les indices $\lfloor n/2\rfloor$ et $\lceil n/2\rceil$, c’est-à-dire :
\begin{itemize}
\item en un point unique si $n$ est impair ;
\item sur un intervalle si $n$ est pair.
\end{itemize}
}

\item   \question{Distinguer les deux cas $n$ impair et $n$ pair pour identifier précisément le ou les minimiseurs de $g$.}
\indication{}
\reponse{
\textbf{Cas 1 : $n=2p+1$ impair.}

On a
\[
\frac n2=p+\frac12.
\]
Sur $(X_{(p)},X_{(p+1)})$, on a
\[
g'(\theta)=2p-(2p+1)=-1<0,
\]
et sur $(X_{(p+1)},X_{(p+2)})$,
\[
g'(\theta)=2(p+1)-(2p+1)=1>0.
\]
Donc $g$ décroît puis croît, et admet un minimum global unique en
\[
\theta=X_{(p+1)}.
\]

\textbf{Cas 2 : $n=2p$ pair.}

Sur $(X_{(p-1)},X_{(p)})$,
\[
g'(\theta)=2(p-1)-2p=-2<0.
\]
Sur $(X_{(p)},X_{(p+1)})$,
\[
g'(\theta)=2p-2p=0,
\]
donc $g$ est constante.
Enfin, sur $(X_{(p+1)},X_{(p+2)})$,
\[
g'(\theta)=2(p+1)-2p=2>0.
\]
Ainsi, $g$ est minimale sur tout le segment
\[
[X_{(p)},X_{(p+1)}].
\]
}

\item   \question{Rappeler la définition de la médiane empirique $\mathrm{med}(X_1,\dots,X_n)$ et conclure que l’estimateur du maximum de vraisemblance est
\[
\widehat{\theta}_{MV}=\mathrm{med}(X_1,\dots,X_n).
\]}
\indication{}
\reponse{
La médiane empirique est définie par
\[
\mathrm{med}(X_1,\dots,X_n)=
\begin{cases}
X_{(p+1)} & \text{si } n=2p+1,\\[1ex]
\dfrac{X_{(p)}+X_{(p+1)}}{2} & \text{si } n=2p.
\end{cases}
\]

D’après la question précédente :
\begin{itemize}
\item si $n$ est impair, l’unique minimiseur de $g$ est $X_{(p+1)}$, qui est la médiane empirique ;
\item si $n$ est pair, tout point de $[X_{(p)},X_{(p+1)}]$ minimise $g$, et en particulier
\[
\frac{X_{(p)}+X_{(p+1)}}{2}
\]
appartient à cet intervalle.
\end{itemize}

Dans les deux cas, la médiane empirique minimise
\[
g(\theta)=\sum_{i=1}^n |X_i-\theta|,
\]
donc maximise la log-vraisemblance. On conclut que
\[
\widehat{\theta}_{MV}=\mathrm{med}(X_1,\dots,X_n).
\]
}

\end{enumerate}

}